% =====================================================================
%  2bat-ud1-15.tex  —  SESSIÓ 15: Taller de preparació de la prova
%  (sense preàmbul: s'inclou des de main.tex)
%  Criteris: C1, C2, C3, C4. Format SIMULACRE (amb solucions per practicar).
% =====================================================================
\horatitol{Sessió 15 — Taller de preparació de la prova}%
{Un simulacre amb un ítem de cada bloc (C1–C4), amb la mateixa estructura que la prova real.}

\subsection*{Com és la prova}

La prova de la propera sessió té \textbf{quatre blocs}, un per cada criteri. Aquest
és un simulacre amb el mateix format perquè hi arribis amb seguretat. Temps
orientatiu: uns \textbf{45 minuts}. Es valora no només el càlcul correcte, sinó
també \emph{comprovar la compatibilitat}, una \emph{disposició clara} i la
\emph{interpretació del resultat} dins el context.

\begin{idea}
\textbf{Estructura i pes orientatiu}
\begin{itemize}[leftmargin=1.4em, itemsep=2pt]
  \item \textbf{Bloc 1 (C1) — Construir i llegir una matriu.} \;(25\%)
  \item \textbf{Bloc 2 (C2) — Sumar/restar i escalar amb interpretació.} \;(25\%)
  \item \textbf{Bloc 3 (C3) — Producte de matrices.} \;(25\%)
  \item \textbf{Bloc 4 (C4) — Compatibilitat i interpretació.} \;(25\%)
\end{itemize}
\end{idea}

\begin{idea}
\textbf{Formulari-resum (el pots usar durant el simulacre)}
\begin{itemize}[leftmargin=1.4em, itemsep=2pt]
  \item \textbf{Dimensió:} files $\times$ columnes. L'element es localitza per
        (fila, columna).
  \item \textbf{Suma/resta:} casella a casella; cal la \emph{mateixa dimensió}.
  \item \textbf{Escalar:} es multiplica \emph{cada} casella pel nombre (retallada
        del $10\%\rightarrow$ multiplicar per $0{,}9$).
  \item \textbf{Producte:} fila per columna (multiplicar element a element i sumar).
        Compatibilitat: columnes de la 1a $=$ files de la 2a. Resultat:
        $(m\times n)\cdot(n\times p)\rightarrow m\times p$.
\end{itemize}
\end{idea}

% ---------------------------------------------------------------------
\subsection{Bloc 1 (C1) — Construir i llegir}

\begin{exercicis}[Bloc 1]
\begin{exlist}
  \item Un menjador social anota els àpats servits en dos torns (migdia i vespre)
        durant tres dies: dilluns $60$ i $40$; dimarts $58$ i $42$; dimecres $62$ i
        $45$. Organitza les dades en una matriu (files: dies; columnes: torns),
        indica'n la dimensió i digues quants àpats es van servir al vespre de
        dimarts.
\end{exlist}
\end{exercicis}

\subsection{Bloc 2 (C2) — Sumar i escalar}

\begin{exercicis}[Bloc 2]
\begin{exlist}
  \item El pressupost de dos serveis (files) en personal i materials (columnes) és
        $A=\begin{pmatrix} 1500 & 400 \\ 1200 & 350 \end{pmatrix}$ i s'hi afegeix una
        partida extraordinària $B=\begin{pmatrix} 300 & 100 \\ 200 & 150 \end{pmatrix}$.
        Troba la despesa total $A+B$ i, després, aplica-hi una retallada lineal del
        $10\%$. Interpreta què significa l'escalar.
\end{exlist}
\end{exercicis}

\subsection{Bloc 3 (C3) — Producte de matrices}

\begin{exercicis}[Bloc 3]
\begin{exlist}
  \item Un servei atén dos usuaris amb tres tipus de sessió:
        $S=\begin{pmatrix} 3 & 1 & 2 \\ 1 & 4 & 0 \end{pmatrix}$. Les durades de cada
        tipus de sessió són $H=\begin{pmatrix} 2 \\ 1 \\ 3 \end{pmatrix}$ hores.
        Calcula $S\cdot H$ i digues quantes hores genera cada usuari.
\end{exlist}
\end{exercicis}

\subsection{Bloc 4 (C4) — Compatibilitat i interpretació}

\begin{exercicis}[Bloc 4]
\begin{exlist}
  \item Una matriu de centres per serveis és $3\times 2$ i una de serveis per cost
        és $2\times 1$. Es pot fer el producte? Quina dimensió tindria el resultat i
        què representaria? I el producte en l'ordre contrari, es podria fer?
\end{exlist}
\end{exercicis}

% ---------------------------------------------------------------------
\fullsolucions{Sessió 15}
\begin{solucions}
\textbf{Bloc 1 (C1)}
\begin{sollist}
  \item $\begin{pmatrix} 60 & 40 \\ 58 & 42 \\ 62 & 45 \end{pmatrix}$,
        \;dimensió $3\times 2$. Al vespre de dimarts (fila 2, columna 2) es van
        servir $42$ àpats.
\end{sollist}

\textbf{Bloc 2 (C2)}
\begin{sollist}
  \item $A+B=\begin{pmatrix} 1800 & 500 \\ 1400 & 500 \end{pmatrix}$. Amb la
        retallada del $10\%$ (escalar $0{,}9$):
        $0{,}9\cdot(A+B)=\begin{pmatrix} 1620 & 450 \\ 1260 & 450 \end{pmatrix}$.
        L'escalar $0{,}9$ aplica el mateix percentatge (deixar el $90\%$) a tots els
        conceptes alhora.
\end{sollist}

\textbf{Bloc 3 (C3)}
\begin{sollist}
  \item $S\cdot H=\begin{pmatrix} 3\per 2+1\per 1+2\per 3 \\ 1\per 2+4\per 1+0\per 3 \end{pmatrix}
        =\begin{pmatrix} 13 \\ 6 \end{pmatrix}$. L'usuari 1 genera $13$ hores;
        l'usuari 2, $6$ hores.
\end{sollist}

\textbf{Bloc 4 (C4)}
\begin{sollist}
  \item Sí: $3\times 2$ per $2\times 1$ dona un resultat $3\times 1$. Representaria,
        per a cada centre (files), el cost total (única columna). En l'ordre
        contrari, $2\times 1$ per $3\times 2$, \emph{no} es pot fer: les columnes de
        la primera ($1$) no coincideixen amb les files de la segona ($3$).
\end{sollist}

\noindent\textbf{Nota per al professorat.} El simulacre no computa com a nota de
prova; serveix per practicar i orientar el suport previ a l'examen. Convé corregir
amb la rúbrica a la vista, valorant la comprovació de la compatibilitat, la
disposició clara i la interpretació dins el context, no només el resultat numèric.
\end{solucions}
